\documentclass[a4paper,10pt]{article}

\usepackage[utf8]{inputenc}
\usepackage[T1]{fontenc}
\usepackage[brazil]{babel}
\usepackage[top=1.3cm, bottom=1.3cm, left=1.6cm, right=1.6cm]{geometry}
\usepackage{tikz}
\usetikzlibrary{shapes.geometric, arrows.meta, positioning}
\usepackage{tcolorbox}
\tcbuselibrary{skins, breakable}
\usepackage{xcolor}
\usepackage{tabularx}
\usepackage{booktabs}
\usepackage{lmodern}
\usepackage{microtype}
\usepackage{enumitem}
\usepackage{multicol}
\usepackage{parskip}

% ── Paleta Terasite ──────────────────────────────────────────
\definecolor{tsazul}{HTML}{0B3D91}
\definecolor{tsazulmed}{HTML}{1A6BBF}
\definecolor{tsazulclaro}{HTML}{E3EFFB}
\definecolor{tsteal}{HTML}{009BAA}
\definecolor{tstealclaro}{HTML}{DFF5F7}
\definecolor{tsverde}{HTML}{27A35C}
\definecolor{tsverdeclaro}{HTML}{E0F5EB}
\definecolor{tsambar}{HTML}{B56F00}
\definecolor{tsambarclaro}{HTML}{FBF0DC}
\definecolor{tsroxo}{HTML}{6B3FA0}
\definecolor{tsroxoclaro}{HTML}{EFE7F7}
\definecolor{tscinza}{HTML}{4A5560}
\definecolor{tscinzaclaro}{HTML}{F2F4F7}

% ── Caixas ───────────────────────────────────────────────────
\tcbset{
  cxa/.style={enhanced,colback=tsazulclaro,colframe=tsazul,
    colbacktitle=tsazul,coltitle=white,
    arc=4pt,boxrule=.9pt,left=7pt,right=7pt,top=4pt,bottom=4pt,
    fonttitle=\bfseries\small},
  cxt/.style={enhanced,colback=tstealclaro,colframe=tsteal,
    colbacktitle=tsteal,coltitle=white,
    arc=4pt,boxrule=.9pt,left=7pt,right=7pt,top=4pt,bottom=4pt,
    fonttitle=\bfseries\small},
  cxv/.style={enhanced,colback=tsverdeclaro,colframe=tsverde,
    colbacktitle=tsverde,coltitle=white,
    arc=4pt,boxrule=.9pt,left=7pt,right=7pt,top=4pt,bottom=4pt,
    fonttitle=\bfseries\small},
  cxam/.style={enhanced,colback=tsambarclaro,colframe=tsambar,
    colbacktitle=tsambar,coltitle=white,
    arc=4pt,boxrule=.9pt,left=7pt,right=7pt,top=4pt,bottom=4pt,
    fonttitle=\bfseries\small},
  cxr/.style={enhanced,colback=tsroxoclaro,colframe=tsroxo,
    colbacktitle=tsroxo,coltitle=white,
    arc=4pt,boxrule=.9pt,left=7pt,right=7pt,top=4pt,bottom=4pt,
    fonttitle=\bfseries\small},
}

% ── Estilos TikZ ─────────────────────────────────────────────
\tikzset{
  bloco/.style={rectangle,rounded corners=3pt,
    minimum width=4.6cm,minimum height=0.65cm,
    text centered,draw,thick,font=\small\bfseries,inner sep=2pt},
  ba/.style={bloco,fill=tsazulclaro,draw=tsazul,text=tsazul},
  bt/.style={bloco,fill=tstealclaro,draw=tsteal,text=tsteal},
  bv/.style={bloco,fill=tsverdeclaro,draw=tsverde,text=tsverde},
  bam/.style={bloco,fill=tsambarclaro,draw=tsambar,text=tsambar},
  br/.style={bloco,fill=tsroxoclaro,draw=tsroxo,text=tsroxo},
  bc/.style={bloco,fill=tscinzaclaro,draw=tscinza,text=tscinza,
             font=\footnotesize\itshape},
  s/.style={-Stealth,thick,color=tscinza},
  sa/.style={-Stealth,thick,color=tsazul},
}

\begin{document}
\pagestyle{empty}
\setlength{\parskip}{0pt}

%% ════════════════════════════════════════════════════════════
%%  PÁGINA 1
%% ════════════════════════════════════════════════════════════

% ── Cabeçalho ────────────────────────────────────────────────
\begin{tcolorbox}[
  enhanced, arc=5pt, boxrule=0pt,
  left=12pt, right=12pt, top=8pt, bottom=8pt,
  interior style={left color=tsazul, right color=tsteal},
  colframe=tsazul,
]
  \color{white}
  {\large\bfseries Protocolo de Testes de Campo --- Antenas Wi-Fi}\\[2pt]
  {\small Frotall FRITG01LTE \quad|\quad Terasite Tecnologia \quad|\quad Abril 2026}
\end{tcolorbox}

\vspace{4pt}

% ── Contexto ─────────────────────────────────────────────────
\begin{tcolorbox}[cxa, title={Contexto}]
\small
Estamos selecionando a melhor antena Wi-Fi para o gateway Frotall, instalado em
máquinas de construção. O operador conecta o celular para sincronizar dados e
configurar a máquina (\textbf{3 a 10 metros} de distância típica).
Testaremos \textbf{6 antenas} (5 externas + 1 interna original) cruzadas com
\textbf{2 estados da máquina} e \textbf{múltiplos pontos de instalação}.
\end{tcolorbox}

\vspace{4pt}
\begin{center}
{\normalsize\bfseries\color{tsazul} Os três testes}
\end{center}
\vspace{2pt}

% ── Teste 1: WALK ─────────────────────────────────────────────
\begin{tcolorbox}[cxv, title={Teste 1 --- WALK (Varredura dinâmica + TTR)}]
\small
\textbf{O que faz:} mede RSSI a 2 Hz enquanto o engenheiro caminha ao redor da
máquina, registrando também \textbf{histerese de desconexão} e \textbf{TTR
(Time-To-Reconnect)} ao sair e voltar para a área de cobertura.\\[2pt]
\textbf{Para que serve:} mapear cobertura geral, identificar zonas mortas e
caracterizar o comportamento de reconexão Wi-Fi (mais difícil de medir do que a
queda em si --- relevante para a UX de sincronização).\\[2pt]
\textbf{Como:} caminhar devagar e constante ($<$ 1 m/s), celular na mesma altura,
trajeto sistemático. \textbf{Duração:} $\sim$3 min por antena.
\end{tcolorbox}

\vspace{3pt}

% ── Testes 2 e 3 lado a lado ──────────────────────────────────
\begin{multicols}{2}

\begin{tcolorbox}[cxt, title={Teste 2 --- CLOCK (Estático por posição)}]
\small
\textbf{O que faz:} mede RSSI a 1 Hz em pontos fixos marcados pelo operador.
O parser agrupa as amostras por marcador.\\[2pt]
\textbf{Para que serve:} comparar sinal a 3 m, 5 m, 10 m, 15 m e em zonas da
máquina (Cabine, Motor, Frente, Trás).\\[2pt]
\textbf{Como:} parar 5--10 s em cada posição \emph{antes} de marcar --- as
amostras seguintes ficam naquele ponto.\\[2pt]
\textbf{Duração:} $\sim$5 min por antena.
\end{tcolorbox}

\columnbreak

\begin{tcolorbox}[cxam, title={Teste 3 --- BURN (Throughput + PLR)}]
\small
\textbf{O que faz:} mede throughput TCP e perda de pacotes em janelas de 15 s
(3 s ping + 10 s tráfego TCP + 2 s cooldown).\\[2pt]
\textbf{Para que serve:} avaliar o desempenho real na sincronização --- o uso
mais pesado do Wi-Fi no Frotall.\\[2pt]
\textbf{Como:} \texttt{watch.py} no notebook antes de iniciar, ESP32 parado,
repetir 3$\times$.\\[2pt]
\textbf{Duração:} $\sim$5 min por antena.
\end{tcolorbox}

\end{multicols}

\vspace{3pt}

% ── Matriz de Estresse Ambiental ─────────────────────────────
\begin{tcolorbox}[cxr, title={Matriz de Estresse Ambiental --- DES $\times$ RUN}]
\small
Cada um dos 3 testes é executado em \textbf{dois estados da máquina}, isolando
a contribuição RF da contribuição mecânica/elétrica:
\begin{itemize}[leftmargin=*,itemsep=1pt,topsep=2pt]
  \item[\textbullet] \textbf{DES (Motor desligado):} baseline puro de RF --- só
        topografia, multipath e ruído ambiente da obra.
  \item[\textbullet] \textbf{RUN (Motor ligado):} adiciona vibração mecânica
        sobre o conector e \textbf{EMI do alternador/ignição} (banda 2.4 GHz é
        notoriamente vulnerável). A diferença DES$\to$RUN é a assinatura de
        ruído da máquina capturada pelo firmware.
\end{itemize}
\end{tcolorbox}

\vspace{3pt}

% ── Antenas ──────────────────────────────────────────────────
\begin{tcolorbox}[cxa, title={Antenas a testar}]
\small
\begin{center}
\begin{tabular}{cll}
\toprule
\textbf{Código} & \textbf{Tipo} & \textbf{Observação} \\
\midrule
A5.3 & Externa & Melhor na bancada (score 100) \\
A5.1 & Externa & Boa na bancada (score 76,7) \\
A5.2, A4.1, A4.2 & Externas & Medianas na bancada \\
INT  & Interna (original) & \textbf{Baseline obrigatório} --- referência da campanha \\
\bottomrule
\end{tabular}
\end{center}
\vspace{1pt}
{\color{tscinza}\footnotesize O teste da INT mostra matematicamente o ganho de
trocar a antena de fábrica por uma externa. Sem ela, a comparação é cega.}
\end{tcolorbox}

\vspace{3pt}

% ── Tempo + Pontos lado a lado ───────────────────────────────
\begin{multicols}{2}

\begin{tcolorbox}[cxam, title={Tempo estimado}]
\small
\begin{tabular}{lc}
\toprule & \textbf{Aprox.} \\ \midrule
WALK (DES+RUN) por ant. & $\sim$6 min \\
CLOCK (DES+RUN) por ant. & $\sim$10 min \\
BURN (DES+RUN) por ant. & $\sim$10 min \\
\midrule
\textbf{Por antena} & $\sim$26 min \\
\textbf{6 antenas + setup} & \textbf{$\sim$4--5 h} \\
\bottomrule
\end{tabular}
\end{tcolorbox}

\columnbreak

\begin{tcolorbox}[cxt, title={Pontos de instalação}]
\small
A posição do ESP32 na máquina afeta o resultado:
\begin{itemize}[leftmargin=*,itemsep=1pt,topsep=2pt]
  \item[\textbullet] \textbf{CABINE} --- posição padrão atual
  \item[\textbullet] \textbf{MOTOR} --- pior caso de EMI
  \item[\textbullet] \textbf{FRENTE} / \textbf{TRAS} --- variação de obstrução
\end{itemize}
\vspace{1pt}
{\color{tscinza}\footnotesize O parser agrupa por marcador --- basta marcar a
posição antes de iniciar o run.}
\end{tcolorbox}

\end{multicols}

\vspace{3pt}
\begin{center}
{\footnotesize\color{tscinza}
Guilherme Bertanha --- Estágio Eng.\ Mecatrônica ---
Terasite Tecnologia --- Abril 2026}
\end{center}

%% ════════════════════════════════════════════════════════════
%%  PÁGINA 2
%% ════════════════════════════════════════════════════════════
\newpage

% ── Cabeçalho ────────────────────────────────────────────────
\begin{tcolorbox}[
  enhanced, arc=5pt, boxrule=0pt,
  left=12pt, right=12pt, top=7pt, bottom=7pt,
  interior style={left color=tsazul, right color=tsteal},
  colframe=tsazul,
]
  \color{white}
  {\normalsize\bfseries Pipeline de Execução + Instrumentação + Entregáveis}\\[-1pt]
  {\small Frotall FRITG01LTE \quad|\quad Terasite Tecnologia \quad|\quad Abril 2026}
\end{tcolorbox}

\vspace{5pt}

% ── Setup de rede ────────────────────────────────────────────
\begin{tcolorbox}[cxa, title={Setup de rede no campo}]
\small
\begin{center}
\begin{tikzpicture}[node distance=0.9cm]
  \node[ba, minimum width=2.7cm, minimum height=0.55cm] (cel) {Celular (hotspot)};
  \node[bt, minimum width=2.7cm, minimum height=0.55cm,
        below left=0.9cm and 1.1cm of cel] (esp)
    {ESP32 (antena em teste)};
  \node[bv, minimum width=2.7cm, minimum height=0.55cm,
        below right=0.9cm and 1.1cm of cel] (nb)
    {Notebook (watch.py)};
  \draw[Stealth-Stealth, thick, color=tsazul] (cel) -- node[above left,font=\tiny,color=tscinza]{Wi-Fi} (esp);
  \draw[Stealth-Stealth, thick, color=tsazul] (cel) -- node[above right,font=\tiny,color=tscinza]{Wi-Fi} (nb);
  \draw[Stealth-Stealth, thick, dashed, color=tsazul] (esp) -- node[below,font=\tiny,color=tscinza]{mesma rede} (nb);
\end{tikzpicture}
\end{center}
{\footnotesize\color{tscinza}
Celular cria o hotspot. ESP32 e notebook entram nessa rede. Sem internet --- tudo
local. Painel web: \texttt{http://\textless IP-ESP32\textgreater/} no celular.}
\end{tcolorbox}

\vspace{5pt}
\begin{center}
{\normalsize\bfseries\color{tsazul} Pipeline de Execução em Campo}
\end{center}
\vspace{3pt}

% ── Pipeline 2 linhas × 3 colunas ────────────────────────────
\begin{center}
\begin{tikzpicture}[
  fase/.style={rectangle, rounded corners=5pt, draw, thick,
    minimum width=4.4cm, minimum height=1.45cm, text centered,
    font=\footnotesize, inner sep=4pt, align=center},
  setaP/.style={-Stealth, line width=1.3pt, color=tscinza},
  font=\footnotesize,
]

  % Linha 1: F1 → F2 → F3
  \node[fase, fill=tsazulclaro, draw=tsazul] (f1) at (0,0)
    {{\large\textbf{\textcolor{tsazul}{F1 \textbullet{} Setup}}}\\[2pt]
     hotspot do celular\\watch.py no notebook};
  \node[fase, fill=tsroxoclaro, draw=tsroxo] (f2) at (5.2,0)
    {{\large\textbf{\textcolor{tsroxo}{F2 \textbullet{} Baseline}}}\\[2pt]
     Antena INT\\\textit{referência da campanha}};
  \node[fase, fill=tstealclaro, draw=tsteal] (f3) at (10.4,0)
    {{\large\textbf{\textcolor{tsteal}{F3 \textbullet{} Estáticos}}}\\[2pt]
     CLOCK por posição\\3 / 5 / 10 / 15 m};

  % Linha 2: F6 ← F5 ← F4 (ordem invertida pra fazer o U)
  \node[fase, fill=tsverdeclaro, draw=tsverde] (f4) at (10.4,-2.6)
    {{\large\textbf{\textcolor{tsverde}{F4 \textbullet{} Dinâmicos}}}\\[2pt]
     WALK + TTR\\BURN (3$\times$ por ant.)};
  \node[fase, fill=tsambarclaro, draw=tsambar] (f5) at (5.2,-2.6)
    {{\large\textbf{\textcolor{tsambar}{F5 \textbullet{} Motor ON}}}\\[2pt]
     repete F3+F4 com\\vibração + EMI 2.4 GHz};
  \node[fase, fill=tsazulclaro, draw=tsazul] (f6) at (0,-2.6)
    {{\large\textbf{\textcolor{tsazul}{F6 \textbullet{} Parser}}}\\[2pt]
     sync.py + relatório\\HTML com veredicto};

  % Setas linha 1
  \draw[setaP] (f1) -- (f2);
  \draw[setaP] (f2) -- (f3);
  % Conexão F3 → F4 (desce)
  \draw[setaP] (f3) -- (f4);
  % Setas linha 2 (direita pra esquerda)
  \draw[setaP] (f4) -- (f5);
  \draw[setaP] (f5) -- (f6);

  % Loop "repetir por antena" — contorna por baixo e pela direita
  \draw[setaP, dashed, color=tsambar, rounded corners=8pt]
    (f5.south) -- (5.2,-4.1) -| (13.4,0) -- (f3.east);

  \node[below=2pt, font=\scriptsize\bfseries, color=tsambar]
    at (7.8,-4.1) {repetir F3$\to$F4$\to$F5 para cada antena externa (5$\times$)};

\end{tikzpicture}
\end{center}

\vspace{2pt}
{\footnotesize\color{tscinza}
A INT (\textbf{F2}) estabelece a referência da campanha. Cada antena externa
passa por \textbf{F3$\to$F4$\to$F5} antes que a próxima entre em teste ---
minimiza trocas físicas, que são o gargalo operacional. O parser (\textbf{F6})
roda offline após tudo coletado.}

\vspace{5pt}

% ── Instrumentação auxiliar (SA6) ────────────────────────────
\begin{tcolorbox}[cxr, title={Instrumentação auxiliar --- Analisador SA6 (35--6200 MHz)}]
\small
O analisador de espectro \textbf{Arinst SSA-TG R2 (SA6)} fica em \emph{standby}
durante a campanha. Acionado para varrer a banda 2.4 GHz \textbf{caso a transição
DES$\to$RUN apresente anomalias severas} (jitter de RSSI, PLR elevado, queda
abrupta de throughput). Captura a \textbf{assinatura de EMI} da máquina real,
distinguindo se a degradação vem de multipath ou de ruído elétrico do motor ---
informação que o firmware sozinho não consegue separar.
\end{tcolorbox}

\vspace{5pt}

% ── Veredicto final dual ─────────────────────────────────────
\begin{tcolorbox}[cxa, title={Entregáveis do relatório final --- Veredicto Dual}]
\small
O parser gera um HTML com banner de vencedor, radar 3 eixos, heatmap antena
$\times$ condição, cobertura por distância, time series de RSSI e ranking. Mas
a recomendação ao time de produto vai em \textbf{dois resultados separados}:

\vspace{3pt}
\begin{tabular}{@{}p{0.46\linewidth}@{\hspace{6pt}}p{0.50\linewidth}@{}}
\textbf{\color{tsteal}1.\ Vencedor de RF (Score Matemático)} &
\textbf{\color{tsverde}2.\ Recomendação de Engenharia} \\[2pt]
{\footnotesize Saída direta do parser:
PLR 35\% + TTR 10\% + RSSI 15\% + Throughput 40\%.
Comparação fria, reprodutível e auditável.} &
{\footnotesize Aplica peso subjetivo sobre robustez do chicote, proteção contra
vibração e viabilidade mecânica observada durante o BURN. Pode divergir do
vencedor matemático.} \\
\end{tabular}

\end{tcolorbox}

\vspace{4pt}
\begin{center}
{\footnotesize\color{tscinza}
Guilherme Bertanha --- Estágio Eng.\ Mecatrônica ---
Terasite Tecnologia --- Abril 2026}
\end{center}

\end{document}
